\documentclass[14pt]{jsarticle}
\input{lecture-preamble.tex}

% 例題・練習の見出し: \reidai{番号}，\renshu{番号}
\newcommand{\reidai}[1]{\par\noindent\textbf{例題\,#1}\quad}
\newcommand{\renshu}[1]{\par\noindent\textbf{練習\,#1}\quad}

% 指数関数の特徴の図: \expfeature{底}{p}{q}{a>1 なら 1，0<a<1 なら 0}（0<a<1 では a^p，a^q を y 軸の右に置く）
\newcommand{\expfeature}[4]{%
  \begin{tikzpicture}[x=0.5cm, y=0.42cm]
    \draw[-{Stealth}] (-3,0) -- (3,0) node[right]{$x$};
    \draw[-{Stealth}] (0,-0.5) -- (0,4.6) node[above]{$y$};
    \node[below left, font=\footnotesize] at (0,0) {O};
    \ifnum#4=1
      \draw[thick, domain=-3:2.15, samples=60] plot (\x, {#1^(\x)});
    \else
      \draw[thick, domain=-2.15:3, samples=60] plot (\x, {#1^(\x)});
    \fi
    \ifnum#4=1 \def\lblside{left}\else\def\lblside{right}\fi
    \draw[dashed] (#2,0) node[below, font=\footnotesize]{$p$} -- (#2,{#1^(#2)}) -- (0,{#1^(#2)})
      node[\lblside, font=\footnotesize]{$a^p$};
    \draw[dashed] (#3,0) node[below, font=\footnotesize]{$q$} -- (#3,{#1^(#3)}) -- (0,{#1^(#3)})
      node[\lblside, font=\footnotesize]{$a^q$};
    \node[above \lblside, inner sep=1pt, font=\footnotesize] at (0,1) {$1$};
  \end{tikzpicture}%
}

\begin{document}

\lectureheader{数学Ⅱ}{5}{指数関数と対数関数}{1}{指数関数}{2-BC}{指数関数の特徴，方程式・不等式}
  {158,159}{指数関数の特徴を使って，大小比較や方程式・不等式について考えよう}

% 同じ節の2枚目以降は，前回までの番号を引き継ぐ
\setcounter{definition}{2}
\setcounter{theorem}{1}
\setcounter{technique}{1}

% ============ 1ページ目（表・左） ============
\heading{〜前時の復習〜}
\begin{theorembox}{指数関数 $y=a^x$ のグラフの特徴}
  \noindent
  \begin{minipage}[t]{0.48\linewidth}
    \textbf{(i)\enspace\boldmath$a>1$ のとき}\par
    \centering
    \begin{tikzpicture}[x=0.5cm, y=0.38cm]
      \draw[-{Stealth}] (-3,0) -- (2.6,0) node[right]{$x$};
      \draw[-{Stealth}] (0,-0.6) -- (0,4.4) node[above]{$y$};
      \node[below left, font=\footnotesize] at (0,0) {O};
      \ifshowanswer
        \begin{scope}[blue]
          \draw[thick, domain=-3:1.75, samples=60] plot (\x, {2.2^(\x)});
          \draw[dashed] (1,0) node[below, font=\footnotesize]{$1$} -- (1,2.2) -- (0,2.2) node[left, font=\footnotesize]{$a$};
          \node[above left, inner sep=1pt, font=\footnotesize] at (0,1) {$1$};
        \end{scope}
      \fi
    \end{tikzpicture}
  \end{minipage}\hfill
  \begin{minipage}[t]{0.48\linewidth}
    \textbf{(ii)\enspace\boldmath$0<a<1$ のとき}\par
    \centering
    \begin{tikzpicture}[x=0.5cm, y=0.38cm]
      \draw[-{Stealth}] (-2.6,0) -- (3,0) node[right]{$x$};
      \draw[-{Stealth}] (0,-0.6) -- (0,4.4) node[above]{$y$};
      \node[below left, font=\footnotesize] at (0,0) {O};
      \ifshowanswer
        \begin{scope}[blue]
          \draw[thick, domain=-1.75:3, samples=60] plot (\x, {0.45^(\x)});
          \draw[dashed] (1,0) node[below, font=\footnotesize]{$1$} -- (1,0.45) -- (0,0.45) node[left, font=\footnotesize]{$a$};
          \node[above right, inner sep=1pt, font=\footnotesize] at (0,1) {$1$};
        \end{scope}
      \fi
    \end{tikzpicture}
  \end{minipage}
  \par\vspace{0.5em}
  \begin{enumerate}[label=\textbf{\arabic*}, leftmargin=1.5zw, itemsep=0.2em]
    \item 2点 $(0,\ \fitblankbf{1})$，$(1,\ \fitblankbf{a})$ を通る．
    \item $\fitblankbf{x \text{軸}}$ を漸近線としてもつ．
    \item $a>1$ のとき\fitblankbf{\text{右上がり}}の曲線，\\
          $0<a<1$ のとき\fitblankbf{\text{右下がり}}の曲線である．
  \end{enumerate}
\end{theorembox}

% \heading{復習問題}
% \textbf{1}\quad 次の数を $2$ の累乗の形で表せ．
% \begin{tasks}[label=(\arabic*), label-width=1.5zw, item-indent=2zw, after-item-skip=0.4em](2)
%   \task $8 = \answermath{2^3}$
%   \task $\dfrac{1}{4} = \answermath{\frac{1}{2^2} = 2^{-2}}$
%   \task $\sqrt{2} = \answermath{2^{\frac{1}{2}}}$
%   \task $\sqrt[3]{4} = \answermath{\sqrt[3]{2^2} = 2^{\frac{2}{3}}}$
% \end{tasks}

% \textbf{2}\quad 次の数を $3$ の累乗の形で表せ．
% \begin{tasks}[label=(\arabic*), label-width=1.5zw, item-indent=2zw, after-item-skip=0.4em](2)
%   \task $\dfrac{1}{9} = \answermath{\frac{1}{3^2} = 3^{-2}}$
%   \task $9^x = \answermath{(3^2)^x = 3^{2x}}$
% \end{tasks}

\vfill
\nextpage

% ============ 2ページ目（表・右） ============
\heading{指数関数の特徴}
\begin{definitionbox}{増加関数・減少関数}
  $x$ の値が増加すると $y$ の値も増加する関数を\fitblankbf{\text{増加関数}}といい，\\
  $x$ の値が増加すると $y$ の値は減少する関数を\fitblankbf{\text{減少関数}}という．
\end{definitionbox}

左のグラフから，$y=a^x$ が増加関数になるか減少関数になるかを考えてみよう．
\par\vspace{0.5em}
\hspace*{1zw}$a>1$ のとき，$y=a^x$ は\fitblankbf{\text{増加関数}}である．\par\vspace{0.6em}
\hspace*{1zw}$0<a<1$ のとき，$y=a^x$ は\fitblankbf{\text{減少関数}}である．
\par\vspace{0.5em}
\begin{theorembox}{指数関数 $y=a^x$ の特徴}
  \begin{enumerate}[label=\textbf{\arabic*}, leftmargin=1.5zw, itemsep=0.3em]
    \item 定義域は実数全体，値域は正の数全体である．
    \item $a>1$ のとき，\fitblankbf{\text{増加関数}}である．\\
          すなわち\quad $p<q \iff a^p \fitblankbf{<} a^q$
    \item $0<a<1$ のとき，\fitblankbf{\text{減少関数}}である．\\
          すなわち\quad $p<q \iff a^p \fitblankbf{>} a^q$
  \end{enumerate}
  \vspace{0.3em}
  \noindent
  \begin{minipage}[t]{0.48\linewidth}
    \textbf{\boldmath$a>1$ のとき}\par\centering
    \expfeature{2}{1}{2}{1}
  \end{minipage}\hfill
  \begin{minipage}[t]{0.48\linewidth}
    \textbf{\boldmath$0<a<1$ のとき}\par\centering
    \expfeature{0.5}{-2}{-1}{0}
  \end{minipage}
\end{theorembox}

% \textbf{注意}\quad $a>0$，$a \neq 1$ のとき，次が成り立つ．\quad
% $p=q \iff a^p \fitblankbf{=} a^q$

% \vspace{0.8em}
% \reidai{2} 次の3つの数の大小を不等号を用いて表せ．
% \hfill $\sqrt{2}$，\ $\sqrt[3]{4}$，\ $\sqrt[5]{8}$\hspace*{3zw}
% \par\vspace{0.8em}
% \textbf{解}\quad 2の累乗で表すと\par\vspace{0.5em}
% \hspace*{3zw}$\sqrt{2} = \answermath{2^{\frac{1}{2}}}$，\
%               $\sqrt[3]{4} = \answermath{2^{\frac{2}{3}}}$，\
%               $\sqrt[5]{8} = \answermath{2^{\frac{3}{5}}}$
% \par\vspace{1em}
% \hspace*{2zw}\begin{tabular}[t]{@{}l@{\hspace{1zw}}l@{}}
%   指数の大小を調べると & $\answermath{\frac{1}{2} < \frac{3}{5} < \frac{2}{3}}$ \\[0.8em]
%   底 $2$ は $1$ より\fitblankbf{\text{大きい}}から & $\answermath{2^{\frac{1}{2}} < 2^{\frac{3}{5}} < 2^{\frac{2}{3}}}$ \\[0.8em]
%   すなわち & $\answermath{\sqrt{2} < \sqrt[5]{8} < \sqrt[3]{4}}$
% \end{tabular}

\begin{techniquebox}{数の大小比較}
  \begin{itemize}[leftmargin=1.5zw, itemsep=0.2em]
    \item 底を\fitblankbf{\text{そろえて}}，指数の大小を比べる．
    \item 底が \fitblankbf{1} より大きいか小さいかで，大小の向きが決まる．
  \end{itemize}
\end{techniquebox}
\vfill
\nextpage

% ============ 3ページ目（裏・左） ============
\heading{指数関数を含む方程式について考えてみよう！}
指数方程式 $8^x = 4$ を解くにはどうしたら良いだろうか．

そもそも、「方程式を解く」ということは、\fitblankbf{x = \text{○}} の形まで持っていくことである．

指数を含む数の大小関係を考えるために，\fitblankbf{\text{底を揃えて}}みよう！
\vspace{0.4em}

\begin{answerproof}[解答]
  方程式を変形すると
  \begin{answeralign}[1zw]
    (2^3)^x &= 2^2 \\
    2^{3x}  &= 2^2 \\
    3x      &= 2
  \end{answeralign}
  よって $x = \fitblankbf{\dfrac{2}{3}}$
\end{answerproof}

となり、指数関数を含む方程式を解くことができた！
\vspace{1em}

\heading{指数関数を含む不等式について考えてみよう！}
$2^x \geqq 8$ を解くにはどうしたら良いだろうか．

まずは、方程式と同様に考えてみよう.

\begin{answerproof}[解答]
  不等式を変形すると
  \begin{answeralign}[1zw]
    2^x &\geqq 2^3
  \end{answeralign}
  \uwave{底 $2$ は $1$ より大きいから}，指数を比べて

  $x \geqq \fitblankbf{3}$
\end{answerproof}

となり、指数関数を含む不等式を解くことができた！しかし、指数を含む不等式を解く際には注意が必要である．

\vfill
\nextpage

% ============ 4ページ目（裏・右） ============


\vspace{0.5em}
\begin{techniquebox}{指数を含む方程式・不等式}
  \begin{itemize}[leftmargin=1.5zw, itemsep=0.2em]
    \item 両辺の\fitblankbf{\text{底}}をそろえて，指数を比べる．
    \item 不等式では，底が $0<a<1$ のとき不等号の向きが\fitblankbf{\text{逆}}になる．
  \end{itemize}
\end{techniquebox}
\vspace{0.3em}

\begin{summarybox}
  \begin{enumerate}[label=\textbf{\arabic*}, leftmargin=1.5zw, itemsep=0.2em]
    \item $a>1$ のとき $y=a^x$ は\fitblankbf{\text{増加関数}}，\\
          $0<a<1$ のとき\fitblankbf{\text{減少関数}}である．
    \item 数の大小比較や方程式・不等式は，\fitblankbf{\text{底}}をそろえて指数を比べる．
  \end{enumerate}
\end{summarybox}


\vfill
\nexttime{対数とは何かを学び，指数との関係を考えてみよう！}

\end{document}
